% ======================================================================
\chapter[Бета-распад ядер. Условия бета-распада]%
        {Бета-распад атомных ядер. Условия бета-распада}
\label{t:20}
% ======================================================================

\begin{qbox}
\textbf{Билет 20.} Бета-распад атомных ядер ($\beta^-$, $\beta^+$, $\mathrm E$-захват), условия бета-распада.
\end{qbox}

\noindent{\small «$\mathrm E$-захват» — электронный захват (захват ядром электрона атомной оболочки; чаще всего с K-оболочки, поэтому
используют и термин «K-захват»).}

\section{Содержание билета}

\textbf{Подтемы:}
\begin{enumerate}
 \item три типа бета-распада: $\beta^-$, $\beta^+$, электронный захват; их нуклонные и кварковые схемы;
 \item энергетические условия бета-распада через массы атомов; распад свободного нейтрона; почему свободный протон стабилен;
 \item непрерывный спектр электронов и гипотеза нейтрино; форма спектра; масса нейтрино; регистрация нейтрино;
 \item электронный захват: условия, сопровождающие явления, случаи «только захват»;
 \item изобарные параболы: бета-стабильность, ядра, распадающиеся обоими способами; времена жизни и правила отбора.
\end{enumerate}

\textbf{Необходимые понятия:} энергия связи и формула Вайцзеккера (\tref{t:17}); слабое взаимодействие (\tref{t:16});
закон распада (\tref{t:18}); фазовый объём (\tref{t:04}); плотность $s$-электронов в ядре (\tref{t:10}).

\section{Определения и физические принципы}

\begin{definition}[бета-распад]
Превращение ядра в изобар с зарядом, отличающимся на единицу, обусловленное слабым взаимодействием:
\begin{align}
  \beta^-:&\quad \nuc{X}{A}{Z}\to\nuc{Y}{A}{Z+1}+e^-+\bar\nu_e &&(n\to p+e^-+\bar\nu_e;\ d\to u+e^-+\bar\nu_e),\label{eq:20:bm}\\
  \beta^+:&\quad \nuc{X}{A}{Z}\to\nuc{Y}{A}{Z-1}+e^++\nu_e &&(p\to n+e^++\nu_e\ \text{в ядре}),\label{eq:20:bp}\\
  \text{E-захват}:&\quad e^-+\nuc{X}{A}{Z}\to\nuc{Y}{A}{Z-1}+\nu_e &&(p+e^-\to n+\nu_e).\label{eq:20:ec}
\end{align}
Массовое число $A$ не меняется; сохраняются электрический заряд, барионное и лептонное числа.
\end{definition}

\begin{definition}[энергия бета-распада]
$Q$ — суммарная кинетическая энергия продуктов, выраженная через массы \emph{атомов} $M(A,Z)$ (см. вывод~\ref{der:20:cond}):
\begin{align}
  Q_{\beta^-}&=\bigl[M(A,Z)-M(A,Z+1)\bigr]c^2, \label{eq:20:Qm}\\
  Q_{\beta^+}&=\bigl[M(A,Z)-M(A,Z-1)-2\mel\bigr]c^2, \label{eq:20:Qp}\\
  Q_{\text{EC}}&=\bigl[M(A,Z)-M(A,Z-1)\bigr]c^2-\varepsilon_i, \label{eq:20:Qec}
\end{align}
где $\varepsilon_i$ — энергия связи захватываемого электрона в оболочке $i$ (K, L,\,\dots). Распад возможен при $Q>0$.
\end{definition}

\emph{Физический смысл.} Бета-распад — способ, которым ядро с «неоптимальным» соотношением $N/Z$ приближается ко дну долины
стабильности (\tref{t:17}): нейтроноизбыточные ядра испытывают $\beta^-$-распад, протоноизбыточные — $\beta^+$-распад или захват.
Электрон и нейтрино не существуют в ядре заранее (электрон в ядре имел бы энергию $\gtrsim40$~МэВ, \tref{t:02}), а рождаются в
момент распада, как фотон при излучении атома.

\emph{Область применимости.} Энергетические условия — точные (с пренебрежением малыми разностями энергий связи атомных
электронов). Форма спектра в простейшей статистической теории — для разрешённых переходов и без учёта кулоновского
взаимодействия электрона с ядром (функции Ферми).

\section{Основная теория}

\subsection{Энергетические условия}

Условия $Q>0$ (вывод~\ref{der:20:cond}) в терминах масс атомов:
\begin{equation}\label{eq:20:conds}
\begin{aligned}
  \beta^-&:\quad M(A,Z)>M(A,Z+1);\\
  \beta^+&:\quad M(A,Z)>M(A,Z-1)+2\mel;\\
  \text{E-захват}&:\quad M(A,Z)>M(A,Z-1)+\varepsilon_i/c^2 .
\end{aligned}
\end{equation}
\textbf{Следствия.}
\begin{enumerate}
 \item Если $\beta^+$-распад возможен, то возможен и захват (условие для захвата слабее на $2\mel c^2=1{,}022$~МэВ); обратное неверно.
 \item При $0<[M(A,Z)-M(A,Z-1)]c^2<1{,}022$~МэВ возможен \emph{только электронный захват}. Пример: $\nucl{Be}{7}\to\nucl{Li}{7}$
       ($Q_{\text{EC}}=0{,}86$~МэВ, $T_{1/2}=53$~сут); $\nucl{Fe}{55}\to\nucl{Mn}{55}$ ($0{,}23$~МэВ).
 \item \textbf{Свободный нейтрон} распадается: $Q=(m_n-m_p-\mel)c^2=0{,}782$~МэВ, $\tau_n\approx880$~с (вывод~\ref{der:20:neutron}).
       \textbf{Свободный протон} стабилен: $m_p<m_n+\mel$. В ядре же протон может превратиться в нейтрон, если энергия связи
       дочернего ядра достаточно велика, т.\,е. энергия берётся из перестройки ядра. Аналогично нейтрон в стабильном ядре не
       распадается, если $M(A,Z+1)>M(A,Z)$.
\end{enumerate}

\subsection{Непрерывный спектр и нейтрино}

\paragraph{Проблема.} $\alpha$-частицы и $\gamma$-кванты моноэнергетичны, а электроны бета-распада имеют \emph{непрерывный} спектр от
нуля до максимальной энергии $T_{\max}\approx Q$ (рис.~\ref{fig:20:spectrum}). Начальное и конечное состояния ядер дискретны, поэтому
двухчастичный распад $X\to Y+e$ давал бы фиксированную энергию электрона. Калориметрические измерения (Эллис и Вустер,
1927) показали, что средняя энергия, выделяемая на один распад, равна среднему $\bar T_e\approx Q/3$, а не $Q$, — часть энергии
«исчезает». Кроме того, нарушался бы закон сохранения момента: например, $\nucl{C}{14}$ и $\nucl{N}{14}$ (чётное $A$) имеют целый спин, а
электрон — полуцелый.

\paragraph{Гипотеза Паули (1930).} Вместе с электроном испускается нейтральная, очень слабо взаимодействующая частица со спином
$\frac12$ и очень малой массой — нейтрино (название дал Ферми). Энергия $Q$ распределяется между тремя телами (электрон,
антинейтрино, ядро отдачи); поскольку ядро тяжёлое, оно почти не уносит энергии, и $T_e+E_\nu\approx Q$. Спектр электронов
непрерывен, а верхняя граница $T_{\max}=Q$ соответствует нулевой энергии нейтрино. Закон сохранения момента
выполняется: $\frac12+\frac12$ — целое.

\paragraph{Теория Ферми (1934).} Распад — превращение $n\to p$ с одновременным рождением пары $e^-\bar\nu$ в одной точке (контактное
слабое взаимодействие, константа $G_F$; на современном языке — обмен виртуальным $W^-$, \tref{t:16}). Вероятность распада
определяется «золотым правилом» $\dd\lambda=\frac{2\pi}{\hbar}|M_{fi}|^2\dd\rho$, где $\dd\rho$ — плотность конечных состояний. Для
разрешённых переходов $|M_{fi}|^2$ от энергии не зависит, и форма спектра определяется фазовым объёмом (вывод~\ref{der:20:spectrum}):
\begin{equation}\label{eq:20:spec}
  \frac{\dd N}{\dd T_e}\propto F(Z,E_e)\,p_eE_e\,(Q-T_e)^2 ,
\end{equation}
где $E_e=T_e+\mel c^2$, $p_e$ — импульс электрона, $F(Z,E_e)$ — функция Ферми, учитывающая притяжение электрона (отталкивание
позитрона) ядром: спектры $\beta^-$ сдвинуты к малым энергиям, $\beta^+$ — к большим.

\begin{figure}[htbp]
\centering
\begin{tikzpicture}
\begin{axis}[width=0.66\linewidth, height=5.0cm, xlabel={$T_e/Q$}, ylabel={$\dd N/\dd T_e$ (отн. ед.)}, xmin=0, xmax=1.05,
   ymin=0, ymax=1.12, samples=200, ytick=\empty, xtick={0,0.25,0.5,0.75,1}]
  \addplot[very thick, accent, domain=0:1] {sqrt(x^2+2*x*0.511)*(x+0.511)*(1-x)^2/0.2523};
  \draw[dashed, gray] (axis cs:1,0) -- (axis cs:1,0.4);
  \node[font=\small, anchor=south] at (axis cs:1,0.4) {$T_{\max}=Q$};
\end{axis}
\end{tikzpicture}
\caption{Форма спектра электронов разрешённого бета-распада по~\eqref{eq:20:spec} без кулоновской поправки
($F=1$), $Q=1$~МэВ. Спектр непрерывен и обрывается при $T_e=Q$; средняя энергия в этом расчёте $\approx0{,}4Q$ (кулоновское притяжение ядра сдвигает спектр $\beta^-$ к меньшим энергиям, и на опыте типично $\bar T_e\approx Q/3$).}
\label{fig:20:spectrum}
\end{figure}

\paragraph{Масса нейтрино.} Если $m_\nu\neq0$, то вблизи верхней границы $(Q-T_e)^2$ заменяется на $(Q-T_e)\sqrt{(Q-T_e)^2-m_\nu^2c^4}$, и
спектр обрывается при $T_e=Q-m_\nu c^2$. Исследование конца спектра трития ($\nucl{H}{3}\to\nucl{He}{3}+e^-+\bar\nu_e$, $Q=18{,}6$~кэВ —
малое $Q$ повышает долю событий у границы) даёт ограничение $m_\nu c^2<0{,}45$~эВ (эксперимент KATRIN). Для удобства анализа
строят \emph{график Кюри}: $\sqrt{N(T)/(pEF)}$ против $T$ — для $m_\nu=0$ и разрешённого перехода это прямая, пересекающая ось при $T=Q$.

\paragraph{Регистрация нейтрино.} Из-за малого сечения ($\sigma\sim10^{-43}$~см$^2$ при энергии $\sim$МэВ) нейтрино прямо
обнаружили лишь в 1956~г. (Ф.~Райнес, К.~Коуэн) у ядерного реактора — мощного источника антинейтрино — по реакции обратного
бета-распада $\bar\nu_e+p\to n+e^+$ (регистрировались совпадения сигнала аннигиляции позитрона и захвата нейтрона).

\subsection{Электронный захват}

Ядро захватывает электрон атомной оболочки, преимущественно с K-оболочки: плотность $s$-электронов в ядре максимальна
и растёт как $Z^3$ (\tref{t:10}). Продукты — ядро $(A,Z-1)$ и моноэнергетическое нейтрино (двухчастичное конечное состояние,
$E_\nu\approx Q_{\text{EC}}$). Сам процесс почти «невидим»; регистрируются сопровождающие явления: на место ушедшего K-электрона
переходят электроны внешних оболочек с испусканием характеристического рентгеновского излучения \emph{дочернего} атома или
оже-электронов. Захват конкурирует с $\beta^+$-распадом: для тяжёлых ядер и малых $Q$ он преобладает (кулоновский барьер
подавляет вылет позитрона, а плотность K-электронов велика). Так как скорость захвата пропорциональна плотности электронов
в ядре, она слегка (порядка $0{,}1\,\%$) зависит от химического окружения — единственный распространённый случай заметной
зависимости $\lambda$ от внешних условий.

\subsection{Изобарные параболы. Бета-стабильность}

Массы изобаров ($A$ фиксировано) по формуле Вайцзеккера — квадратичная функция $Z$ (\tref{t:17}):
\begin{equation}\label{eq:20:parab}
  M(A,Z)c^2=M_0(A)c^2+\frac{\beta(A)}{2}(Z-Z_0)^2-\delta,\qquad \beta(A)=\frac{8a_A}{A}+\frac{2a_C}{A^{1/3}}.
\end{equation}
\begin{itemize}
 \item \textbf{Нечётное $A$} ($\delta=0$): одна парабола. Ядра на левой ветви испытывают $\beta^-$-распад, на правой — $\beta^+$/захват,
       пока не достигнут минимума. Как правило, существует ровно \emph{один} бета-стабильный изобар.
 \item \textbf{Чётное $A$}: две параболы, сдвинутые на $2\delta$ — нижняя для чётно-чётных, верхняя для нечётно-нечётных ядер
       (рис.~\ref{fig:20:parabola}). Распады идут зигзагом между ними. Возможны два и даже три стабильных чётно-чётных изобара
       (их взаимное превращение требует двойного бета-распада, крайне медленного: $T\gtrsim10^{19}$ лет). Нечётно-нечётные ядра
       почти все нестабильны и часто распадаются \emph{обоими} способами: $\nucl{Cu}{64}$ ($\beta^-$ в $\nucl{Zn}{64}$ и $\beta^+$/захват в $\nucl{Ni}{64}$),
       $\nucl{K}{40}$ ($\beta^-$ в $\nucl{Ca}{40}$ — около 89\,\%, захват в $\nucl{Ar}{40}$ — около 11\,\%).
\end{itemize}

\begin{figure}[htbp]
\centering
\begin{tikzpicture}
\begin{axis}[width=0.48\linewidth, height=5.2cm, xlabel={$Z$}, ylabel={$M(A,Z)$}, xmin=-3.6, xmax=3.6, ymin=-0.5, ymax=6.3,
   xtick={-3,...,3}, xticklabels={$Z_0{-}3$,,,$Z_0$,,,$Z_0{+}3$}, ytick=\empty, samples=100,
   title={\small нечётное $A$}]
  \addplot[gray, dashed, domain=-3.4:3.4] {0.55*(x-0.2)^2};
  \addplot[only marks, mark=*, mark size=2.3pt, accent] coordinates {(-3,5.632) (-2,2.662) (-1,0.792) (0,0.022) (1,0.352) (2,1.782) (3,4.312)};
  \draw[->, keycol, thick] (axis cs:-2.9,5.4) -- (axis cs:-2.1,2.9);
  \draw[->, keycol, thick] (axis cs:-1.9,2.45) -- (axis cs:-1.1,1.0);
  \draw[->, keycol, thick] (axis cs:-0.9,0.65) -- (axis cs:-0.1,0.15);
  \draw[->, dercol, thick] (axis cs:2.9,4.1) -- (axis cs:2.1,2.0);
  \draw[->, dercol, thick] (axis cs:1.9,1.6) -- (axis cs:1.1,0.55);
  \draw[->, dercol, thick] (axis cs:0.9,0.3) -- (axis cs:0.1,0.08);
  \node[font=\scriptsize, keycol] at (axis cs:-2.2,4.4) {$\beta^-$};
  \node[font=\scriptsize, dercol] at (axis cs:2.9,3.0) {$\beta^+$, EC};
\end{axis}
\end{tikzpicture}\hfill
\begin{tikzpicture}
\begin{axis}[width=0.48\linewidth, height=5.2cm, xlabel={$Z$}, ylabel={$M(A,Z)$}, xmin=-3.6, xmax=3.6, ymin=-1.3, ymax=6.3,
   xtick={-3,...,3}, xticklabels={,,,,,,}, ytick=\empty, samples=100,
   title={\small чётное $A$}]
  \addplot[gray, dashed, domain=-3.4:3.4] {0.5*x^2-0.9};
  \addplot[gray, dotted, domain=-3.4:3.4] {0.5*x^2+0.9};
  \addplot[only marks, mark=*, mark size=2.3pt, accent] coordinates {(-2,1.1) (0,-0.9) (2,1.1)};
  \addplot[only marks, mark=o, mark size=2.3pt, keycol, thick] coordinates {(-3,5.4) (-1,1.4) (1,1.4) (3,5.4)};
  \draw[->, thick] (axis cs:-0.95,1.25) -- (axis cs:-0.1,-0.65);
  \draw[->, thick] (axis cs:-1.05,1.25) -- (axis cs:-1.85,1.25);
  \node[font=\scriptsize, align=center] at (axis cs:-1.2,3.2) {нечётно-нечётное:\\распад в обе стороны};
  \node[font=\scriptsize, accent] at (axis cs:2.3,-0.6) {чётно-чётные};
  \draw[<->] (axis cs:3.25,3.6) -- (axis cs:3.25,5.4) node[midway, left, font=\scriptsize] {$2\delta$};
\end{axis}
\end{tikzpicture}
\caption{Изобарные параболы (схема). Слева: нечётное $A$ — единственный стабильный изобар у дна параболы.
Справа: чётное $A$ — две параболы (чётно-чётные ядра — заполненные кружки, нечётно-нечётные — пустые), сдвинутые на $2\delta$;
нечётно-нечётное ядро может распадаться и с увеличением, и с уменьшением $Z$; стабильными могут оказаться несколько
чётно-чётных изобаров.}
\label{fig:20:parabola}
\end{figure}

\subsection{Времена жизни и правила отбора}

Периоды бета-распада — от миллисекунд до $10^{16}$ лет и более. Главные факторы: (1) энерговыделение — при $Q\gg\mel c^2$
$\lambda\propto Q^5$ (закон Сарджента, вывод~\ref{der:20:sargent}); (2) правила отбора. Электрон и нейтрино в разрешённых переходах
вылетают с $l=0$ (их длины волн $\gg R$), поэтому момент ядра может измениться лишь на их суммарный спин:
\emph{фермиевские} переходы (спины $e$ и $\nu$ антипараллельны) — $\Delta J=0$; \emph{гамов-теллеровские} (спины параллельны) —
$\Delta J=0,\pm1$ (кроме $0\to0$); чётность ядра не меняется. Переходы с $l\ge1$ — \emph{запрещённые}, они медленнее на порядки
(множитель $\sim(pR/\hbar)^{2l}$). Для сравнения переходов используют «приведённое время жизни» $fT_{1/2}$, где $f(Z,Q)$ —
интеграл по спектру (для разрешённых $\lg fT\approx3$--$6$).

\section{Математические выводы}

\begin{derivation}[label=der:20:cond]{энергетические условия через массы атомов}
\emph{Обозначения:} $M_{\text я}(Z)$ — масса ядра, $M(Z)=M_{\text я}(Z)+Z\mel-B_e(Z)/c^2$ — масса атома ($B_e$ — энергия связи электронов;
разностью $B_e(Z)-B_e(Z\pm1)$ — десятки эВ--кэВ — пренебрегаем). Масса нейтрино $\approx0$.

\emph{$\beta^-$:} $Q=[M_{\text я}(Z)-M_{\text я}(Z+1)-\mel]c^2$. Подставим $M_{\text я}(Z)=M(Z)-Z\mel$, $M_{\text я}(Z+1)=M(Z+1)-(Z+1)\mel$:
\[
  Q_{\beta^-}=\bigl[M(Z)-Z\mel-M(Z+1)+(Z+1)\mel-\mel\bigr]c^2=\bigl[M(Z)-M(Z+1)\bigr]c^2 .
\]
Масса испущенного электрона «компенсируется» лишним электроном в оболочке дочернего атома.

\emph{$\beta^+$:} $Q=[M_{\text я}(Z)-M_{\text я}(Z-1)-\mel]c^2$:
\[
  Q_{\beta^+}=\bigl[M(Z)-Z\mel-M(Z-1)+(Z-1)\mel-\mel\bigr]c^2=\bigl[M(Z)-M(Z-1)-2\mel\bigr]c^2 .
\]
Один $\mel$ — масса позитрона, второй — лишний электрон оболочки, который «освобождается» (дочерний атом имеет на один
электрон меньше нужного).

\emph{Захват:} начальное состояние — ядро $Z$ и один из его электронов; $Q=[M_{\text я}(Z)+\mel-M_{\text я}(Z-1)]c^2-\varepsilon_i$:
\[
  Q_{\text{EC}}=\bigl[M(Z)-Z\mel+\mel-M(Z-1)+(Z-1)\mel\bigr]c^2-\varepsilon_i=\bigl[M(Z)-M(Z-1)\bigr]c^2-\varepsilon_i .
\]
Таблицы содержат именно атомные массы, поэтому формулы~\eqref{eq:20:Qm}--\eqref{eq:20:Qec} удобны на практике.
\end{derivation}

\begin{derivation}[label=der:20:neutron]{распад нейтрона и стабильность протона}
$Q_n=(m_n-m_p-\mel)c^2=939{,}565-938{,}272-0{,}511=0{,}782$~МэВ $>0$ — нейтрон распадается.
Максимальная кинетическая энергия электрона $\approx0{,}782$~МэВ (отдача протона: при максимальном импульсе электрона
$pc=\sqrt{T(T+2\mel c^2)}\approx1{,}19$~МэВ протон получает $p^2/2m_p\approx0{,}75$~кэВ — пренебрежимо).

Для свободного протона: $Q=(m_p-m_n-\mel)c^2=-1{,}80$~МэВ $<0$ — распад $p\to ne^+\nu$ невозможен. Захват $p+e^-\to n+\nu$ (атом водорода):
$Q=(m_{\mathrm H}-m_n)c^2\approx-0{,}78$~МэВ — тоже невозможен. В ядрах эти процессы идут за счёт разности энергий связи:
$Q_{\beta^+}=Q_{\text{нукл}}+[E_{\text{св}}(Z-1)-E_{\text{св}}(Z)]$.
\end{derivation}

\begin{derivation}[label=der:20:spectrum]{форма бета-спектра из фазового объёма}
\emph{Допущения:} разрешённый переход ($|M_{fi}|$ не зависит от энергий); нейтрино безмассово; отдачей ядра пренебрегаем
(оно уносит импульс, но не энергию); кулоновское поле ядра не учитывается ($F=1$).

\emph{Шаг 1.} Число состояний электрона и нейтрино в объёме $V$ (\tref{t:04}):
$\dd n=\dfrac{V\,4\pi p_e^2\dd p_e}{(2\pi\hbar)^3}\cdot\dfrac{V\,4\pi p_\nu^2\dd p_\nu}{(2\pi\hbar)^3}$ (импульс ядра фиксирован законом сохранения).

\emph{Шаг 2.} Энергия сохраняется: $T_e+E_\nu=Q$, $E_\nu=p_\nu c$. При фиксированном $p_e$ плотность состояний на единицу полной
энергии получается заменой $\dd p_\nu=\dd E_\nu/c$: $\dd\rho\propto p_e^2\dd p_e\cdot p_\nu^2/c$, где $p_\nu=(Q-T_e)/c$.

\emph{Шаг 3.} Переход к переменной $T_e$: из $E_e^2=p_e^2c^2+\mel^2c^4$ $\Rightarrow$ $p_e\dd p_e=E_e\dd E_e/c^2=E_e\dd T_e/c^2$. Итого
\[
  \frac{\dd N}{\dd T_e}\propto p_e^2\frac{\dd p_e}{\dd T_e}(Q-T_e)^2=p_eE_e(Q-T_e)^2/c^2 .
\]
\emph{Анализ:} при $T_e\to0$ $\dd N/\dd T_e\propto p_e\to0$ (мало состояний электрона); при $T_e\to Q$ $\propto(Q-T_e)^2\to0$ (мало состояний
нейтрино). Максимум — в средней части спектра; при $Q=1$~МэВ и $F=1$ численное интегрирование даёт $\bar T_e\approx0{,}4Q$, а с учётом функции Ферми для $\beta^-$-распада — около $Q/3$.
\end{derivation}

\begin{derivation}[label=der:20:sargent]{закон Сарджента для бета-распада}
Полная вероятность $\lambda\propto\int_0^Q p_eE_e(Q-T_e)^2\dd T_e$. В ультрарелятивистском пределе $Q\gg\mel c^2$: $p_ec\approx E_e\approx T_e$:
\[
  \lambda\propto\int_0^QT^2(Q-T)^2\dd T=Q^5\int_0^1x^2(1-x)^2\dd x=Q^5\Bigl(\frac13-\frac24+\frac15\Bigr)=\frac{Q^5}{30}.
\]
Таким образом $\lambda\propto G_F^2Q^5$ — результат, согласующийся с размерным анализом для распада мюона (\tref{t:16}).
Удвоение энергии распада уменьшает время жизни в $\approx30$ раз. При $Q\lesssim\mel c^2$ электрон нерелятивистский, зависимость от $Q$ другая, и закон $Q^5$ даёт лишь грубую оценку.
\end{derivation}

\begin{derivation}[label=der:20:parab]{бета-стабильность и изобарная парабола}
Из~\eqref{eq:20:parab} при нечётном $A$ разность масс соседних изобаров
\[
  [M(Z)-M(Z+1)]c^2=\frac\beta2\bigl[(Z-Z_0)^2-(Z+1-Z_0)^2\bigr]=\beta\Bigl(Z_0-Z-\frac12\Bigr).
\]
$\beta^-$-распад $Z\to Z+1$ возможен при $Z<Z_0-\frac12$; аналогично $\beta^+$/захват $Z\to Z-1$ — при $Z>Z_0+\frac12$ (для захвата;
для $\beta^+$ — при большем удалении, с учётом $2\mel c^2$). Стабилен изобар с $|Z-Z_0|<\frac12$ — как правило, единственный.

При чётном $A$ для нечётно-нечётного ядра $Z$ оба соседа — чётно-чётные и лежат на параболе ниже на $2\delta$:
$[M(Z)-M(Z\pm1)]c^2=2\delta+\beta(\dots)$. Если $2\delta$ превышает разность ординат параболы, распад энергетически возможен в обе
стороны — так распадаются $\nucl{Cu}{64}$ и $\nucl{K}{40}$.
\end{derivation}

\section{Итоговые формулы}

\keyformula{Q_{\beta^-}=[M(Z)-M(Z+1)]c^2,\quad Q_{\beta^+}=[M(Z)-M(Z-1)-2\mel]c^2,\quad Q_{\text{EC}}=[M(Z)-M(Z-1)]c^2-\varepsilon_i}
 {$M$ — массы нейтральных атомов, $\varepsilon_i$ — энергия связи захватываемого электрона.}
 {пренебрегаем разностью энергий связи атомных электронов и массой нейтрино.}
 {энергетические условия и энерговыделение всех типов бета-распада.}

\keyformula{\frac{\dd N}{\dd T_e}\propto F(Z,E_e)\,p_eE_e(Q-T_e)^2,\qquad T_{\max}=Q-m_\nu c^2}
 {$T_e$, $E_e$, $p_e$ — кинетическая, полная энергия и импульс электрона; $F$ — функция Ферми.}
 {разрешённые переходы; отдача ядра пренебрежима.}
 {форма непрерывного спектра; масса нейтрино по концу спектра.}

\keyformula{n\to p+e^-+\bar\nu_e:\ Q=0{,}782\ \text{МэВ},\ \ \tau\approx880\ \text{с};\qquad \lambda\propto Q^5\ (Q\gg\mel c^2)}
 {$Q$ — энергия распада.}
 {закон Сарджента — для разрешённых переходов при больших $Q$.}
 {стабильность нуклонов; оценка относительных времён жизни.}

\keyformula{M(A,Z)c^2=M_0c^2+\frac\beta2(Z-Z_0)^2-\delta,\qquad Z_0\approx\frac{A}{1{,}98+0{,}015A^{2/3}}}
 {$\beta=8a_A/A+2a_C/A^{1/3}$, $\delta$ — энергия спаривания.}
 {капельная модель.}
 {какие изобары стабильны, направление бета-распада, двусторонние распады нечётно-нечётных ядер.}

\section{Физическая интерпретация}

\textbf{Роль бета-распада.} Сильное взаимодействие формирует ядро, но не может менять $N/Z$; это делает слабое взаимодействие.
Поэтому бета-распады медленные (от миллисекунд до миллиардов лет), а распределение стабильных ядер на плоскости $(N,Z)$ —
узкая долина, по склонам которой ядра «скатываются» бета-распадами.

\textbf{Нейтрино.} Непрерывный спектр — главный экспериментальный аргумент в пользу нейтрино. Оно уносит в среднем около
двух третей энергии распада, что важно для энергетики реакторов и звёзд (нейтринные потери).

\textbf{Сравнение трёх процессов.} $\beta^-$ — избыток нейтронов; $\beta^+$ и захват — избыток протонов, причём при $\Delta Mc^2<1{,}022$~МэВ
возможен только захват. Захват даёт моноэнергетические нейтрино и рентгеновское излучение дочернего атома; $\beta^+$ —
аннигиляционные $\gamma$-кванты по $0{,}511$~МэВ (основа позитронно-эмиссионной томографии).

\section{Что необходимо знать на экзамене}

\begin{examcheck}
\checkhead{Определения}
\begin{itemize}
 \item схемы $\beta^-$, $\beta^+$, E-захвата на уровне ядер, нуклонов и кварков; законы сохранения ($Q$, $B$, $L_e$, момент);
 \item непрерывный спектр, гипотеза нейтрино (энергия и момент), обнаружение нейтрино (Райнес--Коуэн).
\end{itemize}
\checkhead{Формулы}
\begin{itemize}
 \item условия распада через атомные массы: $M(Z)>M(Z+1)$; $M(Z)>M(Z-1)+2\mel$; $M(Z)>M(Z-1)+\varepsilon/c^2$;
 \item $Q_n=0{,}782$~МэВ; форма спектра $p_eE_e(Q-T_e)^2$; $\lambda\propto Q^5$; изобарная парабола.
\end{itemize}
\checkhead{Выводы}
\begin{itemize}
 \item переход от ядерных масс к атомным в условиях $\beta^\pm$ и захвата (откуда $2\mel$);
 \item форма спектра из фазового объёма; закон Сарджента (интеграл $Q^5/30$);
 \item бета-стабильность для нечётных и чётных $A$.
\end{itemize}
\checkhead{Интерпретации}
\begin{itemize}
 \item почему свободный протон стабилен, а в ядре возможен $\beta^+$-распад; случаи «только захват» ($\nucl{Be}{7}$);
 \item двусторонние распады ($\nucl{Cu}{64}$, $\nucl{K}{40}$); правила отбора Ферми и Гамова--Теллера; масса нейтрино по спектру трития.
\end{itemize}
\end{examcheck}
